\section*{Chapitre \ref{ch:ml:neural-networks-for-noisy-tabular-data} --- Réseaux de neurones pour des données tabulaires bruitées}

\begin{solution}{exo:ml:neural-networks-for-noisy-tabular-data:1}
$40\times32 + 32 + 32\times16 + 16 + 16 + 1 = 1\,857$. Avec une première couche de 64~: $40\times64 + 64 + 64\times32 + 32 +
32\times16 + 16 + 17 = 5\,249$.
\end{solution}

\begin{solution}{exo:ml:neural-networks-for-noisy-tabular-data:2}
$n \ge 8\times0.003^2/0.001^2 = 72$ graines par architecture.
\end{solution}

\begin{solution}{exo:ml:neural-networks-for-noisy-tabular-data:3}
$160\,000/512 = 312.5$, soit 313 pas~; les $200\times500 = 100\,000$ lignes du chapitre donnent 196 pas par \omterm{def:ml:neural-networks-for-noisy-tabular-data:backprop}{époque}.
\end{solution}

\begin{solution}{exo:ml:neural-networks-for-noisy-tabular-data:4}
L'IC ne dépend que du classement des prévisions~; le R-carré dépend aussi de leur échelle. Des graines différentes s'arrêtent à
des \omterm{def:ml:neural-networks-for-noisy-tabular-data:backprop}{époques} différentes, avec des prévisions plus ou moins dispersées autour d'un classement semblable, et l'erreur d'échelle entre dans le R-carré
quadratiquement (la \cref{prop:ml:why:ceiling}).
\end{solution}

\begin{solution}{exo:ml:neural-networks-for-noisy-tabular-data:5}
La \omterm{def:ml:neural-networks-for-noisy-tabular-data:dropout}{normalisation par lot} centre et réduit chaque unité avec les statistiques propres au lot, qui sont des estimations bruitées sur 512 lignes
composées surtout de bruit, et elle couple les lignes d'un même lot~; un réseau peu profond sur des entrées classées ne gagne rien
à un meilleur conditionnement en contrepartie. Dans un réseau profond, la normalisation stabilise les échelles des couches et rend le plus souvent
l'entraînement possible tout court~; son bruit est alors un prix qui vaut la peine d'être payé.
\end{solution}

\begin{solution}{exo:ml:neural-networks-for-noisy-tabular-data:6}
La meilleure des vingt scores de validation est un bruit sélectionné (la \cref{prop:ml:validation:max})~: la graine retenue ne fait pas mieux
hors échantillon que les autres. Il faut déployer l'ensemble à vingt graines, qui, lui, l'est.
\end{solution}

\begin{solution}{exo:ml:neural-networks-for-noisy-tabular-data:7}
\texttt{ensemble(n)}~: une graine 0.0602, deux 0.0625, cinq 0.0659, dix 0.0662. Deux graines permettent d'obtenir $(0.0625 - 0.0602)/(0.0662 -
0.0602) = 38~\%$ du gain, cinq graines 95~\%.
\end{solution}

\begin{solution}{exo:ml:neural-networks-for-noisy-tabular-data:8}
Avec $z_1 = \mathrm{relu}(A_1x + c_1)$ et la prévision linéaire en $z_1$, la règle de la chaîne donne $\partial\ell_i/\partial A_1 =
(\partial\ell_i/\partial z_{1,i}\odot\mathbf 1\{A_1x_i + c_1 > 0\})\,x_i^\top$, la dérivée de la ReLU étant la
fonction indicatrice. La passe arrière est un produit matrice-vecteur par couche, comme la passe avant, plus des produits terme à terme~:
un petit multiple constant du coût de la passe avant, indépendant du nombre de paramètres (livre~4, chapitre~28).
\end{solution}

\begin{solution}{pb:ml:neural-networks-for-noisy-tabular-data:1}
\begin{enumerate}
\item 1\,857 et 5\,249.
\item Après deux à quatre \omterm{def:ml:neural-networks-for-noisy-tabular-data:backprop}{époques}~: le score de validation culmine presque d'emblée sur des données aussi bruitées, et les \omterm{def:ml:neural-networks-for-noisy-tabular-data:backprop}{époques} suivantes ajustent le bruit.
\item Le \omterm{def:ml:neural-networks-for-noisy-tabular-data:dropout}{dropout} 0.2 porte l'IC moyen sur trois graines de 0.059 à 0.062~; la \omterm{def:ml:neural-networks-for-noisy-tabular-data:adam}{décroissance des poids} 0.01 ne change rien~; la \omterm{def:ml:neural-networks-for-noisy-tabular-data:dropout}{normalisation
par lot} l'abaisse à 0.052, la \omterm{def:ml:neural-networks-for-noisy-tabular-data:dropout}{normalisation par couche} à 0.056.
\item L'écart-type des graines (0.0024) est comparable aux écarts, si bien que trois graines donnent des erreurs types
de même taille que les effets.
\item IC de 0.057 à 0.065, R-carré de 0.08~\% à 0.30~\%, ratio de Sharpe net de 1.37 à 1.78.
\item 0.0601 et 0.0619~; écart 0.0019~; écart-type commun 0.0024.
\item 14 par architecture.
\item $0.00244\sqrt2 = 0.0035$, soit presque le double de l'effet.
\item 0.066 chacun (le petit ensemble~: R-carré de 0.32~\%, ratio de Sharpe net de 1.73).
\item La moyenne élimine la part propre à chaque graine de l'erreur de chaque membre (la \cref{prop:ml:trees:bagging}).
\item L'écart disparaît~: les deux ensembles atteignent 0.066.
\item La dispersion des prévisions des membres pour chaque ligne, mesure de l'incertitude propre du modèle (chapitre~11).
\item Un réseau seul égale la régression ridge (IC 0.060)~; l'ensemble égale l'IC du lasso (0.066) et est devancé par les arbres boostés
par défaut (0.075, ratio de Sharpe net 2.02).
\item \emph{Résultat nommé}~: l'écart-type des graines de l'IC de test, 0.0024, dépasse l'écart entre architectures, 0.0019~;
distinguer cet écart à deux erreurs types demande 14 graines par architecture, et un ensemble de dix graines (0.066) bat
les exécutions isolées des deux architectures.
\item Les arbres boostés~; si c'est un réseau, son ensemble de graines, jamais sa meilleure graine.
\item L'architecture, chaque \omterm{def:ml:why-financial-machine-learning-is-different:sets}{hyperparamètre}, les graines, le bloc de validation, la meilleure \omterm{def:ml:neural-networks-for-noisy-tabular-data:backprop}{époque} de chaque graine, les empreintes de
hachage des états et la dispersion des résultats.
\item Davantage de données réduisent la variance de chaque ajustement, si bien que les graines comptent moins et que les réseaux plus grands peuvent se rentabiliser.
\item Que le même code, les mêmes données et la même graine produisent les mêmes poids~; non que le résultat soit bon, ni qu'une autre graine
donnerait le même résultat.
\item Les poids initiaux, les masques de \omterm{def:ml:neural-networks-for-noisy-tabular-data:dropout}{dropout}, l'ordre des \omterm{def:ml:neural-networks-for-noisy-tabular-data:backprop}{mini-lots} (et, sur les accélérateurs, les
noyaux non déterministes).
\item Un tirage dans une distribution de modèles, dont la dispersion doit être mesurée avant que sa moyenne soit comparée.
\end{enumerate}
\end{solution}

\begin{solution}{iq:ml:neural-networks-for-noisy-tabular-data:1}
Le gradient de la perte par rapport à chaque poids, obtenu par la règle de la chaîne appliquée de la sortie vers l'entrée
(différentiation en mode inverse). Son coût est un petit multiple d'une passe avant, car la dérivée de chaque couche est réutilisée par toutes les
couches qui la précèdent.
\iqlookfor{le mode inverse et l'argument sur le coût.}
\end{solution}

\begin{solution}{iq:ml:neural-networks-for-noisy-tabular-data:2}
\omterm{def:ml:neural-networks-for-noisy-tabular-data:adam}{Adam} avec une pénalité $\ell_2$ ajoute $\lambda\theta$ au gradient, lequel est ensuite remis à l'échelle par le pas propre à chaque paramètre d'\omterm{def:ml:neural-networks-for-noisy-tabular-data:adam}{Adam}, si bien que
les poids dont la variance du gradient est grande sont à peine pénalisés~; AdamW réduit directement les poids, en dehors du pas adaptatif,
ce qui rétablit le sens de la \omterm{def:ml:neural-networks-for-noisy-tabular-data:adam}{décroissance des poids}.
\iqlookfor{l'interaction entre la pénalité et la mise à l'échelle adaptative.}
\end{solution}

\begin{solution}{iq:ml:neural-networks-for-noisy-tabular-data:3}
Entraîner plusieurs graines de chaque architecture (la \cref{prop:ml:nets:seeds} dit combien), comparer les moyennes avec l'erreur type de la
différence, et comparer les ensembles de graines~; vérifier le gain sur une autre période.
\iqlookfor{les graines comme bruit et une comparaison correcte sur deux échantillons.}
\end{solution}

\begin{solution}{iq:ml:neural-networks-for-noisy-tabular-data:4}
Fixer la graine de chaque générateur (torch, NumPy, l'échantillonneur de lots), utiliser un seul fil d'exécution et des algorithmes déterministes
(\texttt{torch.use\_deterministic\_algorithms}),
figer les versions des bibliothèques et calculer l'empreinte des poids. Des versions de bibliothèques, du matériel, des nombres de fils ou des implémentations BLAS différentes peuvent
encore modifier les résultats aux derniers bits.
\iqlookfor{toute source d'aléa et de non-associativité des calculs en virgule flottante.}
\end{solution}

\begin{solution}{iq:ml:neural-networks-for-noisy-tabular-data:5}
Les arbres~: des valeurs par défaut solides, un réglage peu coûteux, une robustesse au changement d'échelle, des \omterm{def:ml:trees-and-boosting:monotone}{contraintes de monotonie}, une inférence rapide, et en général au moins aussi bons
sur des variables classées. Les réseaux~: interactions lisses, plongements, entraînement conjoint avec d'autres objectifs, un atout avec
de très grands volumes de données, et des ensembles qui estiment leur propre incertitude. Sur un panel mensuel, commencer par les arbres et ajouter un
ensemble de réseaux s'il améliore le stacking.
\iqlookfor{un argument équilibré lié à la taille des données et à l'usage.}
\end{solution}

\begin{solution}{iq:ml:neural-networks-for-noisy-tabular-data:6}
L'erreur de chaque réseau comporte une part commune à toutes les graines (ce que permettent les données et l'architecture) et une part propre à chaque graine~;
la moyenne élimine la seconde, la variance devenant $\rho\sigma^2 + (1-\rho)\sigma^2/B$. Elle cesse d'aider lorsque la part propre à chaque graine est
disparue, après une poignée de graines ici~: l'erreur restante des membres est corrélée.
\iqlookfor{la décomposition de la variance et son plancher.}
\end{solution}
